\section{Dataset}
\label{sec:dataset}
Our dataset is designed to simulate realistic court decision-making in the Indian legal context, incorporating facts, statutes, and precedent, which are essential elements under the common law framework. This dataset enables exploration of Legal Judgment Prediction (LJP) in a Retrieval-Augmented Generation (RAG) setup.

\subsection{Dataset Compilation}

We curated a large-scale dataset consisting of 56,387 Supreme Court of India (SCI) case documents up to April 2024, sourced from IndianKanoon\footnote{\url{https://indiankanoon.org/}}, a trusted legal search engine. The website provides structural tags for various judgment components (e.g., facts, issues, arguments), which allowed for clean and structured scraping. These documents serve as the foundation for our summarization, retrieval, and reasoning experiments.

\subsection{Single versus Multi Partitioning}

Another challenge in the SCI proceedings is the presence of cases containing multiple petitions under a single case identifier, where each petition may lead to distinct outcomes. To handle this, we partition the dataset into two subsets. The \textit{single} partition represents binary classification tasks (\texttt{accept}/\texttt{reject}) and includes cases with a single petition or multiple petitions that share the same outcome. The \textit{multi} partition
%, a superset of \textit{single},
contains cases with multiple petitions leading to different decisions, introducing a more complex multi-label prediction setting (e.g., partially accepted outcomes). For these cases, if even one appeal within the group is accepted, the case is assigned the label \texttt{accepted}. Predicting mixed outcomes remains challenging and will be further explored in future work.

\subsection{Dataset Composition}

The corpus supports multiple downstream pipelines, each focusing on specific judgment elements or legal context. Table~\ref{tab:test-cases} presents key statistics across different configurations.
An example case is shown in Table~\ref{case-example} in Appendix.

\subsubsection{Case Text}

Each judgment includes complete narrative content such as factual background, party arguments, legal issues, reasoning, and verdict. Due to length constraints exceeding model context windows, we summarized these documents using \texttt{Mixtral-8x7B-Instruct-v0.1} \citep{Jiang2024MixtralOE}, which supports up to 32k tokens. The summarization preserved critical legal elements through carefully designed prompts (see Table~\ref{tab:Instruction-sets-summarize}).

%%%%%%%%%%%%%%%%%%%%%%%%%%%%%%%%%%%%%%%%%%%%%%%%%%%%%%%%%%%%%%%%%%%
\begin{table}[t]
\centering
\resizebox{\columnwidth}{!}{
\begin{tabular}{lrrr}
\hline
\textbf{Dataset}  & \textbf{\#Documents} & \textbf{Avg. Length} & \textbf{Max} \\ \hline
\textbf{SCI (Full)} & 56,387 & 3,495 & 401,985  \\
\textbf{Summarized Single} & 4,962 & 302 & 875  \\
\textbf{Summarized Multi}  & 4,930 & 300 & 879  \\
\textbf{Sections} & 29,858 & 257 & 27,553  \\ \hline
\end{tabular}
}
\vspace*{-2mm}
\caption{NyayaRAG data statistics}
\label{tab:test-cases}
\end{table}
%%%%%%%%%%%%%%%%%%%%%%%%%%%%%%%%%%%%%%%%%%%%%%%%%%%%%%%%%%%%%%%%%%%

%%%%%%%%%%%%%%%%%%%%%%%%%%%%%%%%%%%%%%%%%%%%%%%%%%%%%%%%%%%%%%%%%%%
\begin{table*}[t]
\centering
\begin{tabular}{p{0.90\textwidth}}
\toprule
\textbf{Summarization Prompt} \\
\midrule
\footnotesize
The text is regarding a court judgment for a specific case. Summarize it into 1000 tokens but more than 700 tokens. The summarization should highlight the Facts, Issues, Statutes, Ratio of the decision, Ruling by Present Court (Decision), and a Conclusion. \\
\bottomrule
\end{tabular}
\vspace*{-2mm}
\caption{Instruction prompt used with \texttt{Mixtral-8x7B-Instruct-v0.1} for summarizing legal judgments}
\label{tab:Instruction-sets-summarize}
\end{table*}
%%%%%%%%%%%%%%%%%%%%%%%%%%%%%%%%%%%%%%%%%%%%%%%%%%%%%%%%%%%%%%%%%%%

\subsubsection{Precedents}

From each judgment, cited precedents were extracted using metadata tags provided by IndianKanoon. These citations represent explicit legal reasoning and are retained for use during inference to replicate how courts consider prior judgments.

\subsubsection{Statutes}

Statutory references were also programmatically extracted, including citations to laws like the Indian Penal Code and the Constitution of India. Where statute sections exceeded length limits, they were summarized using the same LLM pipeline. Only statutes directly cited in the respective cases were retained, ensuring relevance.

\subsubsection{Previous Similar Cases}
To simulate implicit precedent-based reasoning, we employed semantic similarity retrieval to identify relevant previous cases beyond explicit citations:

\begin{itemize}
    \item \textbf{Corpus Vectorization:} All 56,387 documents were embedded into dense vector representations using the \texttt{all-MiniLM-L6-v2} sentence transformer.
    \item \textbf{Target Encoding:} The 5,000 selected training samples were vectorized similarly.
    \item \textbf{Top-k Retrieval:} Using \texttt{ChromaDB}, we retrieved the top-3 most semantically similar cases for each document based on cosine similarity.
    \item \textbf{Augmentation:} Retrieved cases were appended to the factual input to form the ``casetext + previous similar cases'' input during model inference.
\end{itemize}

This retrieval step enriches context with precedents that are semantically close; even if not cited, they enhance the legal realism of our setup.

\subsubsection{Facts}
We separately extracted the factual portions of all 56,387 judgments. These include background information, chronological events, and party narratives, excluding legal reasoning. These fact-only subsets were used to simulate realistic courtroom scenarios where judges primarily rely on facts, relevant law, and precedent for decision-making.

% \noindent
Overall, our dataset is uniquely structured to test legal decision-making under realistic constraints, aligning with the Indian legal system's reliance on factual narratives, statutory frameworks, precedents, and prior rulings.

% \section{Dataset}

% %%%%%%%%%%%%%%%%%%%%%%%%%%%%%%%%%%%%%%%%%%%%%%%%%%%%%%%%%%%%%%%%%%%%%%%%%%%%%%%%%%%%%%%%%%%%%%%%%%%%%55

% \begin{table*}[ht]
% \centering
% \begin{tabular}{|l|p{\textwidth}|}
% \hline
% \textbf{Summarization prompt} \\ \hline

% The text is regarding a court judgment for a specific case. Summarize it into 1000 tokens\\but more than 700 tokens. The summarization should highlight the Facts, issues, Statute,\\Ratio of the decision, Ruling by Present Court (Decision) and a Conclusion \\ \hline


% \end{tabular}
% \caption{Instruction set used for \texttt{Mixtral 8x7B Instruct-v0.1} inference across all data pipelines.}
% \label{tab:Instruction-sets-summarize}
% \end{table*}


% %%%%%%%%%%%%%%%%%%%%%%%%%%%%%%%%%%%%%%%%%%%%%%%%%%%%%%%%%%%%%%%%%%%%%%%%%%%%%%%%%%%%%%%%%%%%%


% The data set used in this research was taken from trusted online sources. This comprehensive data set offers unparalleled diversity and coverage, addressing existing gaps and providing a rich foundation for advanced legal AI research.

% \subsection{Dataset Compilation}
% The dataset preparation method entailed the collection of a corpus that included 56,387 Supreme Court case proceedings up to April 2024. We used the IndianKanoon website\footnote{\url{https://indiankanoon.org/}}, a prominent legal search engine for the collection of these documents. This source is renowned for its extensive database of Indian legal documents, making it an indispensable asset for our data set. The IndianKanoon website provided tags to multiple parts of the judgment including facts, issues, arguments, and so on which helped us to scrap the data with clear distinction.

% \subsection{Dataset Composition} 
% The dataset was structured to support multiple pipeline configurations, each focusing on different aspects of legal reasoning and citation context. Various parts of a judgment can be observed in the Appendix table~\ref{case-example}. The primary components of the data set are:

% %%%%%%%%%%%%%%%%%%%%%%%%%%%%%%%%%%%%%%%%%%%%
% % Temporal Data Stats

% \begin{table}[t]
% \centering
% \resizebox{\columnwidth}{!}{
% \begin{tabular}{lrrrr}
% \hline
% \textbf{Dataset}  &
%   \textbf{\#Documents} &
%   \textbf{Average} &
%   \textbf{Maximum} &
%   \textbf{Minimum}
%    \\ \hline
% \textbf{SCI}      & 56,387    & 3495   & 4,01,985   & 14   \\ 
% \textbf{Summarized Single}  & 4,962 & 302 & 875 & 1 \\ 
% \textbf{Summarized Multi} & 4,930   & 300   & 879   & 1   \\ 
% \textbf{Sections}  & 29,858 & 257 & 27,553 & 9 \\ \hline
% \end{tabular}%
% }
% \caption{Statistics of data across different categories for evaluation.}
% \label{tab:temporal-test-cases}
% \end{table}
% %%%%%%%%%%%%%%%%%%%%%%%%%%%%%%%%%%%%%%%%%%%%%

% \subsubsection{Case Text}
% Each legal document contains the complete text of the case, including the factual background, arguments from both parties, legal issues, judicial reasoning, and the final verdict. These documents were normalized to ensure consistency in format, with sections clearly delineated where possible. Since the data had word count beyond the context limit of open-source LLMs, we had to summarize the data using the Mixtral-8x7B-Instruct-v0.1 model with context window of 32,000. The exact prompt used has can be found in table~\ref{tab:judgment_prediction_prompts_few}

% \subsubsection{Cited cases}
% The cited cases are extracted from the judgment texts, identifying references to prior case law. These citations are crucial for understanding the precedential value and legal reasoning in judgments. The website contained a separate section for cited cases, ensuring precedential value. The cases were maintained along with the judgment to be used while inference.


% \subsubsection{Statutes}
% Statutory references within the judgments are identified and extracted to form the statutes component. These statutes include provisions from various laws such as the Indian Penal Code, the Constitution of India, and other relevant legislation. The scrapping process ensures that all statutory references are accurately captured. It must be strictly noted that the statutes are limited to the context of a particular judgment. In some instances, the statute information was over limits, which was then summarized using open source LLM model i.e. Mixtral-8x7B-Instruct-v0.1.

% \subsubsection{Previous Similar Cases}

% To enrich the legal context and simulate precedent-aware legal reasoning, a semantic similarity retrieval mechanism was employed to extract previous similar cases for each document in the dataset. This approach goes beyond explicit citations and leverages latent similarities between cases. The methodology is as follows:

% \begin{itemize}
%     \item \textbf{Vectorization of the Corpus:} All 56,387 documents from the scraped corpus were embedded into dense vector representations using a transformer-based sentence embedding model named all-MiniLM-L6-v2. Each document was represented as a single high-dimensional vector encoding its semantic content.

%     \item \textbf{Target Document Embedding:} Each of the 5,000 selected documents for model training was also vectorized using the same embedding model.

%     \item \textbf{Similarity Search via ChromaDB:} To retrieve semantically similar documents, \texttt{ChromaDB}, a vector-native database optimized for high-performance similarity search, was used. For every target document, a top-$k$ retrieval (with $k=3$) was performed to extract the three most semantically similar cases from the full 56k corpus.

%     \item \textbf{Augmentation:} The retrieved similar cases were appended as contextual input to the model during inference, forming the ``\texttt{casetext + previous similar cases}''.
% \end{itemize}

% This technique allows the system to incorporate implicit precedents that may not be directly cited but share legal and factual proximity, thereby enhancing the depth of legal understanding in downstream tasks.

% \subsubsection{Facts}

% For all 56,387 documents scraped from Indian Kanoon, only the factual portions of each case were independently extracted. These sections typically include the background of the case, the circumstances that led to the dispute, and a summary of the events presented by the parties involved. The extraction was done programmatically by identifying and isolating the factual narrative from the rest of the judgment text. This component provides a structured overview of the circumstances of the case without including judicial reasoning, legal citations, or judgments.






